\scanpage{139}
\mspage{10}
% archive sheet number: 138

Quand $A$ est divisible, \deleted{\unsure{on}} $Z\in\unsure{L|\mathcal{L}}$, et on trouve la forme

\deleted{\unsure{(ii'')|(iii')}}\inserted{(iv)} Il n'existe pas $\unsure{x|\alpha}\in L$, avec $x\in\operatorname{Omb}(F)$ (i.e. $x\ll F$), et $x\mathrel{|\circ|}G$.

En fait, les implications
\[
\text{(i)}\Longrightarrow\text{(ii)}\Longrightarrow
\deleted{\text{\unsure{(ii')}}}\text{(iii)}\Longrightarrow
\deleted{\text{\unsure{(ii'')}}}\text{(iv)}
\]
sont \deleted{\unsure{est}} triviales, et $(\deleted{\unsure{\mathrm{iii}}}\mathrm{iv})\Longrightarrow\text{(i)}$ (moyennant divisibilité) \unsure{resp.} $\text{(iii)}\Longrightarrow\text{(i)}$, grâce à \unsure{Atom.g|Atmag|Atomag} supp$^2$.

Signalons enfin une autre condition équivalente, moyennant \uline{divisibilité} :

\deleted{\unsure{(ii quater)|(ii''')}}\inserted{(v)} $G$ est maximal parmi les raffinements de $F$, ordonnés par $\deleted{\unsure{\ll}}\subset$.
% margin note written sideways
% slanted
\margin{\uline{NB} Ces conditions impliquent déjà que $\unsure{\varphi|\phi}$ est \unsure{surjective}.}

\deleted{Mais}\inserted{Quand} ces conditions \deleted{\unsure{n'impliquent pas}}\inserted{sont satisfaites}, on dira \deleted{\unsure{pas encore}} que $G$ est une \uline{subdivision de} $F$. \deleted{(Comme précédemment)}

% The following entire block is crossed through with diagonal strokes.
\deleted{mais je voudrais la condition}

\deleted{(ii) La top. de $\widetilde{F}$ (associée à sa structure d'ordre) est \uline{quotient} de celle de $\widetilde{G}$, ou encore, si $X,X'\in\widetilde{F}$, pour qu'on ait $X\leq X'$ il faut (et suffit) qu'il existe $Y\in\varphi^{-1}(X)$, $Y'\in\varphi^{-1}(X')$, avec $Y\leq Y'$.}
% margin note written sideways
% slanted
\margin{\unsure{est-elle indépendante} de l'exigence \unsure{?} \uline{Voire !}}

\deleted{Cela signifie aussi, dans le cas $A=\unsure{\operatorname{Fig\acute{e}lem}|\operatorname{Fig\acute{e}l\acute{e}m}}(\unsure{E|L}$}
